\chapter{Construire et calculer les nombres complexes}
\section*{Exercice 25.1 -- Calcul algébrique}
Le produit vaut $2+4i-3i-6i^2=8+i$. Multiplier le quotient par le conjugué du dénominateur donne $(1+i)(2+i)/5=(1+3i)/5$. Le dénominateur réel vaut $|2-i|^2=5$.

\section*{Exercice 25.2 -- Conjugué}
Écrire $z=x+i y$ donne $3x-i y=3-i$. L'identification des parties réelle et imaginaire impose $x=1$, $y=1$. La solution unique est $z=1+i$.

\section*{Exercice 25.3 -- Inverse et module}
$z\overline z=|z|^2>0$ ; multiplier $z$ par le quotient indiqué donne $1$. Puis $1=|zz^{-1}|=|z||z^{-1}|$ et la division par $|z|>0$ donne la formule du module.

\section*{Exercice 25.4 -- Binôme complexe}
On peut élever deux fois au carré : $(1+i)^2=2i$, puis $(2i)^2=-4$. Le binôme donne $1+4i+6i^2+4i^3+i^4=1+4i-6-4i+1=-4$.

\section*{Exercice 25.5 -- Absence d’ordre compatible}
Dans un corps ordonné, tout carré est positif ou nul : si $x\ge0$ cela suit de la multiplication ; sinon $-x>0$ et $x^2=(-x)^2>0$. Or $i^2=-1$ serait positif, alors que $1>0$ impose $-1<0$. Cette contradiction interdit un tel ordre.

\section*{Exercice 25.6 -- Inégalité triangulaire}
$|z+w|^2=|z|^2+|w|^2+2\Re(z\overline w)\le|z|^2+|w|^2+2|z\overline w|=(|z|+|w|)^2$. Les deux membres sont positifs, donc on peut prendre les racines. Pour $z=1,w=i$, $\sqrt2<2$.
